\begin{ExerciseBlock}
\ExerciseHeading{03}{Lineare Regression für Verkaufsdaten}{laboratory_measurements_umbrella_sales_regression}
Die zugehörige Excel-Datei auf Moodle enthält monatliche Niederschlagsmengen und Verkaufszahlen für Regenschirme. Die Tabelle zeigt den in der Vorlage abgebildeten Auszug.
\begin{center}
\begin{tabular}{lrr}
\toprule
Monat & Niederschlag in \unit{\milli\metre} & Verkaufte Regenschirme\\
\midrule
Januar & 82 & 15\\
Februar & 92,5 & 25\\
März & 83,2 & 17\\
April & 97,7 & 28\\
Mai & 131,9 & 41\\
Juni & 141,3 & 47\\
Juli & 165,4 & 50\\
August & 140 & 46\\
September & 126,7 & 37\\
\bottomrule
\end{tabular}
\end{center}
\begin{enumerate}
\item Verwenden Sie den bereitgestellten Python-Code, um die vollständigen Daten zu plotten. Schätzen Sie ein, ob ein linearer Zusammenhang plausibel ist. Bestimmen Sie die Steigung $m$ und den Achsenabschnitt $b$ der Regressionsgeraden und tragen Sie sie zusammen mit den Messpunkten auf.
\item Prognostizieren Sie mit Ihrer Regressionsgeraden die Anzahl verkaufter Regenschirme bei \qty{162}{\milli\metre} Niederschlag.
\end{enumerate}
\end{ExerciseBlock}
\UseExerciseSolution{laboratory_measurements_umbrella_sales_regression}
